%قبل از شروع مسئله ترجیحاً چیزی ننویسید. اما اگر خواستید بنویسید، در قالب لیستی (فقرات و شمارش) بنویسید.
%برای اضافه کردن نمره، بعدی از شروع‌مسئله، [x] را اضافه کنید و به جای x  نمره مورد نظر را قرار دهید. برای زیرمسئله‌ها هم به همین صورت است.
%\شروع‌مسئله[10]
\شروع‌مسئله [40]
%\متن‌سیاه{سوال عملی اول:}\\
%\متن‌سیاه{دراپ‌تست}\\

سروش و علیرضا و مهدی و اسماعیل قصد دارند بازی چهارنفره‌ای را انجام دهند، در این بازی یکی از کارت‌ها بسیار محبوب است و همه قصد‌ دارند این کارت را بردارند برای اینکه در حق کسی بی‌انصافی صورت نگیرد، تصمیم می‌گیرند که برای اخذ این کارت به صورت تصادفی عمل کنند.

متاسفانه چون علیرضا که در اینگونه بازی‌ها آوازه‌ی خوبی ندارد، افراد برای جلوگیری از تقلب،‌ تصمیم می‌گیرند به گونه‌ی عجیبی یک نفر را به تصادف انتخاب کنند. چهار رشته‌ی بزرگ عدد درست کرده و هر کس یکی از آن‌ها را انتخاب کند. فردی که بزرگترین زیررشته‌ی پالیندروم در رشته‌اش موجود باشد، برنده است.

یک رشته را پالیندروم گوییم اگر از دو طرف به یک شکل خوانده شود. در فارسی کلمه‌ی داماد و در انگلیسی کلمه‌ی testset نمونه‌هایی از کلمات پالیندروم هستند.


حال علیرضا که به این سادگی‌ها حاضر نیست بیخیال برنده شدن شود، از شما می‌خواهد کدی بزنید که با گرفتن یک رشته‌ی ورودی از حروف کوچک انگلیسی اندازه‌ی بزرگترین زیررشته‌ی پالیندروم را بگوید.

\متن‌سیاه{ورودی:}\\
یک رشته با طول حداکثر
$5 \times 10^4$
است. اجزای رشته جروف کوچک انگلیسی هستند.


\متن‌سیاه{خروجی:}\\
یک عدد که اندازه‌ی بزرگترین زیررشته‌ی پالیندروم در رشته‌ی ورودی می‌باشد.


\متن‌سیاه{محدودیت‌ها:}\\
استفاده از داده ساختار dictionary پایتون مجاز نمی باشد. برای مقایسه رشته ها اجازه استفاده از عملگر == را ندارید. همچنین اجازه استفاده از داده ساختارهایی که در کلاس درس تدریس نشده اند را ندارید.

\متن‌سیاه{ورودی و خروجی نمونه:}\\

تست اول
\begin{latin}
	\begin{tabularx}{\textwidth}{|X | X|}
		\hline
		\textbf{stdin} &
		\textbf{stdout}\\
		\hline
		pqhlettelm & 6\\
		\hline
	\end{tabularx}
\end{latin}


\solution{

ابتدا به کمک الگوریتم rolling hash پیش پردازشی از $O(n)$ انجام می‌دهیم و مقدادیر هش مربوط به رشته و معکوس آن را نگه می‌داریم. پس از این به بعد می‌توانیم در $O(1)$ هش مربوط به یک رشته و هش معکوس آن را بیابیم. پس می‌توانیم در $O(1)$ بررسی کنیم که آیا هش یک رشته با هش معکوس آن برابر است، و پالیندروم بودن آن را تشخیص بدهیم.

برای یافتن بزرگترین زیررشته‌ی پالیندروم، ابتدا مرکز رشته را تعریف می‌کنیم: اگر طول یک رشته $l$ فرد باشد حرف $\frac{l+1}{2}$ و اگر زوج باشد،‌ دو حرف $\frac{l}{2}$ و $\frac{l}{2}+1$ را مرکز رشته می‌نامیم.

حال توجه کنید که اگز یک زیررشته‌ی پالیندروم داشته باشیم، تمام زیررشته‌های این زیررشته با همان مرکز نیز پالیندروم می‌باشند. لذا برای هر مرکز، با استفاده از جست‌وجوی دو دویی روی طول، بزرگترین زیررشته به آن مرکز را در $O(\log(n))$ بیابیم.

از آنجایی که تعداد مراکز زوج و فرد در بین زیررشته‌های یک رشته از $O(n)$ است، با حالت بندی روی مرکز، بزرگترین زیررشته‌ی به آن مرکز را پیدا کرده و بیشینه میگیریم.
 
کد نمونه:‌
\\

\definecolor{codegreen}{rgb}{0,0.6,0}
\definecolor{codegray}{rgb}{0.5,0.5,0.5}
\definecolor{codepurple}{rgb}{0.58,0,0.82}
\definecolor{backcolour}{rgb}{0.95,0.95,0.92}

\lstdefinestyle{mystyle}{
    backgroundcolor=\color{backcolour},   
    commentstyle=\color{codegreen},
    keywordstyle=\color{magenta},
    numberstyle=\tiny\color{codegray},
    stringstyle=\color{codepurple},
    basicstyle=\ttfamily\footnotesize,
    breakatwhitespace=false,         
    breaklines=true,                 
    captionpos=b,                    
    keepspaces=true,                 
    numbers=left,                    
    numbersep=5pt,                  
    showspaces=false,                
    showstringspaces=false,
    showtabs=false,                  
    tabsize=2
}

\lstset{style=mystyle}

\lr{\lstinputlisting[language=python]{src/palind.py}}


}


%}
% صورت اصلی مسئله را در این‌جا بنویسید.
%اگر مسئله چند قسمت داشت به این صورت عمل کنید:
%\شروع{زیرمسائل}
%\زیرمسئله[0] % مسئلۀ ۱آ
%\زیرمسئله % مسئلۀ ۱ب
%\زیرمسئله % مسئلۀ ۱پ
%\زیرمسئله % مسئلۀ ۱ت
%\پایان{زیرمسائل}



%%%%%%%%%%%%%%%%%%%%%%%%%%%%%%%%%%%%%%%%%%%%%%%%%%%%%
% برخی دستورات مفید برای تمرین‌ها. %
%%%%%%%%%%%%%%%%%%%%%%%%%%%%%%%%%%%%%%%%%%%%%%%%%%%%%

% فرمول‌های نمونه:
%$f(x) = \Theta(x)$
%$T(x, y) \leq \log(x) + 2^y + \binom{2n}{n}$
% {\tt code\_function}


% لیست موارد بدون شماره‌گذاری:
%\شروع{فقرات}
%    \فقره مورد اول 
%        \شروع{فقرات}
%            \فقره مورد اول داخلی 
%                \شروع{فقرات}
%                    \فقره اولین مورد خیلی داخلی 
%                    \فقره دومین مورد خیلی داخلی 
%                \پایان{فقرات}
%            \فقره مورد دوم داخلی 
%        \پایان{فقرات}
%    \فقره مورد دوم 
%\پایان{فقرات}

% لیست موارد شماره‌گذاری شده:
%\شروع{شمارش}
%    \فقره یک مورد 
%    \فقره یک مورد دیگر 
%    \فقره باز هم یک مورد دیگر
%\پایان{شمارش}

% برای نوشتن یک فرمول چند خطی که یک شماره بخوره. برای حالت‌های دیگر به سایت زیر مراجعه کنید:
% https://www.overleaf.com/learn/latex/Aligning_equations_with_amsmath
%\begin{align} 
%    \begin{split}
%        (x+y)^3 &= (x+y)^2(x+y) \\
%                &= (x^2+2xy+y^2)(x+y) \\
%                &= (x^3+2x^2y+xy^2) + (x^2y+2xy^2+y^3) \\
%                &= x^3+3x^2y+3xy^2+y^3
%    \end{split}                                 
%\end{align}

% نوشتن ماتریس و آرایه دوبعدی
%\begin{align}
%    A = 
%    \begin{bmatrix}
%        A_{11} & A_{21} \\
%        A_{21} & A_{22}
%    \end{bmatrix}
%\end{align}

% اضافه کردن تصویر با شرح. برای اضافه کردن تصویر بدون شرح، فقط دستور درج‌تصویر را استفاده کنید.
%\شروع{شکل}[h!]
%\تنظیم‌ازوسط
%  \درج‌تصویر[پهنا=0.5\پهنای‌متن]{img.jpg}
%\شرح{شرح تصویر}
%\برچسب{تصویر:۱}
%\پایان{شکل}
%\پایان{زیرمسائل}

%کد زیر را ببینید.
%%\lstinputlisting{source_filename.py}
%\begin{latin}
%\begin{lstlisting}
%# Import the modules
%import sys
%import random
%
%ans = True
%
%while ans:
%question = raw_input("Ask the magic 8 ball a question: (press enter to quit) ")
%
%answers = random.randint(1,8)
%
%if question == "":
%sys.exit()
%
%elif answers == 1:
%print "It is certain"
%
%elif answers == 2:
%print "Outlook good"
%
%elif answers == 3:
%print "You may rely on it"
%
%elif answers == 4:
%print "Ask again later"
%
%elif answers == 5:
%print "Concentrate and ask again"
%
%elif answers == 6:
%print "Reply hazy, try again"
%
%elif answers == 7:
%print "My reply is no"
%
%elif answers == 8:
%print "My sources say no"
%\end{lstlisting}
%\end{latin}
